\documentclass[11pt]{article}
\usepackage[margin=0.75in]{geometry}
\usepackage{graphicx}
\usepackage{booktabs}
\usepackage{amsmath}
\usepackage{hyperref}
\usepackage{caption}
\usepackage{subcaption}
\usepackage{enumitem}
\usepackage{parskip}
\usepackage{float}

\setlength{\parskip}{4pt}
\setlength{\parindent}{0pt}

\title{\vspace{-1.5em}\textbf{PatchTST: A Time Series is Worth 42 Words}\\}
\author{Sajal Sabat \and Ashir Rao}
\date{CS 4782, Spring 2026}

\begin{document}
\maketitle
\vspace{-1.5em}

\section*{1. Introduction}

The usefulness of Transformers in timeseries forecasting became publicly contested when Zeng et al.\ (DLinear, AAAI 2023) showed that a single-layer linear model, DLinear, outperformed every prior Transformer-based forecaster. They concluded transformers are not useful for time-series forecasting. \textbf{``A Time Series is Worth 64 Words: Long-term Forecasting with Transformers''} (Nie, Nguyen, Sinthong, \& Kalagnanam, ICLR 2023) challenges Zeng's claim, arguing prior Transformers failed due to poor tokenization, not the transformer architecture itself. Their model, \textbf{PatchTST}, fixes this by 1). \textit{patching}, which groups timesteps into overlapping subsequences to reduce sequence length and group semantic information and 2). \textit{channel-independence} which means processing each variable independently. This was found to reduce overfitting. Put together, these innovations restore Transformers as state-of-the-art timeseries forecasters.

\section*{2. Chosen Result}

We reproduce \textbf{Table 2} of Nie et al., which reports MSE for PatchTST, DLinear, and prior Transformer attempts on Electricity and Weather for $T \in \{96, 336\}$, $L=336$. This table was the paper's primary evidence that patching puts Transformers over the DLinear challenge, so it was important to reproduce. We also reproduced the lookback length ablation from \textbf{Figure 4}, which tests the role of context length in prediction MSE. A really big benefit of patching is allowing us to have a longer lookback window because each patch is made from several timesteps, so testing the Transformer's newfound lookback abilities was an important part of why PatchTST does better than other Transformers. 

\section*{3. Replication Methodology}

\textbf{Architecture.} Our primary model, \textbf{PatchTST/42} (CI), uses $P=16$, $S=8$, $L=336$, yielding $N=42$ tokens and 921K parameters. A 3-layer Transformer encoder ($d=128$, 16 heads, FFN 256) with learned positional encodings feeds a linear forecast head. We also replicated the paper's self-supervised pretraining, where the same model is first trained to reconstruct masked patches on unlabeled series, then fine-tuned on the forecasting loss. Finally, we trained three lookback ablation variants with different values of $L$.

\begin{table}[H]
\centering\small
\caption{Model variants. Weather dataset, Adam ($\text{lr}=10^{-4}$), early stopping, normalized MSE/MAE.}
\label{tab:variants}
\begin{tabular}{lccccr}
\toprule
Variant & Purpose & Params & $N$ & Attn Pairs & s/ep \\
\midrule
PatchTST/42 (CI) & Main model & 921K & 42 & 1,764 & 26 \\
SSL + fine-tune & Masked pretraining & 921K & 42 & 1,764 & 30 \\
No-patch Transformer & Ablate patching & 4.58M & 337 & 113,569 & 431 \\
Channel-mixing TST & Ablate channel indep. & 11.3M & 42 & 1,764 & 28 \\
DLinear & Paper baseline & {--} & {--} & {--} & {--} \\
\bottomrule
\end{tabular}
\end{table}

\textbf{Datasets \& Metrics.} Electricity (321 variates) and Weather (21 variates), split 70 train/10 val/20 test, normalized per-variate on train statistics. We report MSE/MAE on the test split in z-score units, averaged over 3 seeds (21, 42, 67). Per-cell spread is $\leq 0.002$ MSE, below the $\sim$0.05 gap to DLinear.

\textbf{Modifications from the paper.} We used $d=128$ (vs.\ paper's $d=512$) and a 20-epoch SSL pretraining schedule, far shorter than the paper's due to compute limits. 

\section*{4. Results \& Analysis}

\textbf{Main results.} Table~\ref{tab:main} compares PatchTST/42 against DLinear, a naive baseline, and the paper. We match or beat the paper on 3 of 4 settings. The largest gap is Electricity $T=96$ (0.136 vs.\ 0.130), attributable to our smaller model ($d=128$ vs.\ $d=512$). Our PatchTST beats DLinear on all settings, confirming the paper's central claim.

\begin{table}[H]
\centering
\small
\caption{Test MSE ($L=336$). Paper values from Nie et al.\ Table 2.}
\label{tab:main}
\begin{tabular}{llcccc}
\toprule
Dataset & $T$ & Naive & DLinear & Ours & Paper \\
\midrule
Electricity & 96  & 1.608 & 0.142 & \textbf{0.136} & 0.130 \\
Electricity & 336 & 1.630 & 0.169 & \textbf{0.166} & 0.167 \\
Weather     & 96  & 0.257 & 0.143 & \textbf{0.148} & 0.152 \\
Weather     & 336 & 0.374 & 0.251 & 0.253 & 0.249 \\
\bottomrule
\end{tabular}
\end{table}

\begin{figure}[H]
  \centering
  \begin{subfigure}[t]{0.48\linewidth}
    \includegraphics[width=\linewidth]{required images/electricity_main_results.png}
    \caption{Electricity results vs.\ paper.}
  \end{subfigure}\hfill
  \begin{subfigure}[t]{0.48\linewidth}
    \includegraphics[width=\linewidth]{required images/weather_main_results.png}
    \caption{Weather results vs.\ paper.}
  \end{subfigure}
\end{figure}

\textbf{Lookback ablation.} MSE improves monotonically with longer lookback ($L=96 \to 192 \to 336$: $0.173 \to 0.154 \to 0.149$ on Weather $T=96$), matching Nie et al.'s Figure 4. Patching makes the longer lookback affordable. The no-patch Transformer needs 431 s/epoch vs.\ 26 s for PatchTST/42 and lands at worse MSE.

\begin{figure}[H]
  \centering
  \includegraphics[width=0.6\linewidth]{required images/ablation.png}
  \caption{Lookback length ablation on Weather ($T=96$).}
  \label{fig:ablation}
\end{figure}

\textbf{Self-supervised pretraining.} A 20-epoch masked-patch pretrain followed by fine-tuning reached MSE 0.156 on Weather $T=96$ vs.\ 0.148 cold-starto

\textbf{Extensions.} We decided to try \textbf{Patch shuffling} to test whether the encoder uses inter-patch order or just treats patches as a bag. We permuted a fraction of patch tokens at inference, and found this harms MSE. MSE climbs 0.148 $\to$ 0.269 (+82\%) from 0\% to 100\% shuffle. We also tested different \textbf{Channel groupings} to try to mix the separate channel data  We tried 4 configurations: cluster the channels by Pearson correlation into $3$ groups, $7$ groups, mix all channels, and stay independent across all. Chart 3b shows that any mixing harms MSE, confirming the paper's so far untested claim that channel mixing causes overfitting.

\begin{figure}[H]
  \centering
  \begin{subfigure}[t]{0.48\linewidth}
    \includegraphics[width=\linewidth]{required images/patch_shuffle_ablation.png}
    \caption{Patch shuffle ablation.}
  \end{subfigure}\hfill
  \begin{subfigure}[t]{0.48\linewidth}
    \includegraphics[width=\linewidth]{required images/channel_mixing_ablation.png}
    \caption{Channel grouping ablation.}
  \end{subfigure}
  \caption{Extension experiments on Weather ($T=96$, $L=336$).}
  \label{fig:exts}
\end{figure}

\section*{5. Reflections}

\textbf{Key results.} Results fall within 0.006 MSE of the paper on all four settings, confirmed across using 3 seeds. Extending lookback from $L=96$ to $L=336$ lowers MSE from 0.173 to 0.149 on Weather $T=96$ (a 14\% reduction). With regards to extensions, shuffling patch order raises MSE by 82\% (0.148 $\to$ 0.269) and mixing channels raises MSE by 20\% (0.148 $\to$ 0.178).

\textbf{Lessons learned.}  We learned that patching helps with both accuracy and efficiency. The Transformer with no patching required 431 s/epoch vs.\ 26 s for PatchTST/42 \emph{and} ended at higher MSE, a cost difference we did not anticipate before running the ablation. We also learned that channel independence is important for training transformer models on timeseries, without this we overfit. 


\textbf{Future directions.}  \textit{Longer horizons}: evaluating $T=720$ would reveal how/if PatchTST degrades at extremely far horizons and whether patching helps there. \textit{Adaptive patch sizing per channel}: different variates have different periodicity, and letting patch size $P$ be learned per channel could improve both accuracy and interpretability. \textit{Extended SSL pretraining}: the 0.156 vs.\ 0.152 gap suggests meaningful gains with a larger pretrain budget.

\section*{References}
\begin{enumerate}[leftmargin=1.5em, topsep=2pt, itemsep=1pt]
  \item Nie, Y., Nguyen, N. H., Sinthong, P., \& Kalagnanam, J. (2023). A Time Series is Worth 64 Words: Long-term Forecasting with Transformers. \textit{ICLR 2023}.
  \item Zeng, A., Chen, M., Zhang, L., \& Xu, Q. (2023). Are Transformers Effective for Time Series Forecasting? \textit{AAAI 2023}.
  \item Wu, H., Xu, J., Wang, J., \& Long, M. (2021). Autoformer: Decomposition Transformers with Auto-Correlation for Long-Term Series Forecasting. \textit{NeurIPS 2021}.
  \item Zhou, H. et al.\ (2021). Informer: Beyond Efficient Transformer for Long Sequence Time-Series Forecasting. \textit{AAAI 2021}.
  \item Paszke, A. et al.\ (2019). PyTorch: An Imperative Style, High-Performance Deep Learning Library. \textit{NeurIPS 2019}.
  \item Dua, D. \& Graff, C. (2019). UCI Machine Learning Repository. University of California, Irvine. \textit{Electricity dataset}.
  \item Wetterstation Jena. Max Planck Institute for Biogeochemistry. \textit{Weather dataset} (jena\_climate\_2009\_2016.csv).
\end{enumerate}

\end{document}
